
\chapter{Phase III Completion}

This part reports the work completed after the mid-term review. It supersedes the
mid-term statement wherever the two differ, and each section says what changed.

\section{System architecture and design}

The mid-term architecture is unchanged: cached decoding, one distributed training
engine, separate evaluation paths, and a registry that gates every reported number.

Four components were added. \path{scripts/swap_localisation_test.py} runs the
intervention experiment that measures what a prediction depends on.
\path{scripts/build_composite_cache.py} and \path{scripts/make_leaf_masks.py}
construct the background-replacement augmentation. \path{scripts/build_thesis.py}
and \path{scripts/make_thesis_figures.py} generate the final report and its
figures from the same artifacts. The registry grew from a table renderer into the
document of record, now backing 23 individual claims.

The design principle held: nothing reports a number it cannot trace to a file.

\section{Methodology and implementation}

Three methodological additions were made.

\textbf{Intervention design.} Asking what a model depends on is normally done by
inspecting saliency maps, which requires human judgement about where a heat map
points. We replaced that with an intervention. Each image receives a donor of a
different class twice, once its leaf and once its background. Adoption means the
prediction moves to the donor's class. A same-class donor measures the disruption
caused by the edit itself, and chance adoption is $1/38$, about 2.6\,\%.

\textbf{Pre-registration.} Hypothesis H3 and the result that would falsify it were
recorded before the augmentation experiment ran, so the prediction could not be fitted
to the outcome afterwards.

\textbf{Paired augmentation design.} Six runs, three seeds per condition, one frozen
protocol, differing only in whether background replacement is enabled. The unit of
independence for the paired interval is the seed.

\section{Prototype development}

The prototype is unchanged from mid-term and the statements there remain accurate: the
research pipeline is complete, the static front end runs, and it shows recorded
checkpoint output rather than live inference because the serving layer is not built.

Added in Phase III: the front end's demo mode is now backed by assets generated by
\path{scripts/make_demo_assets.py} from a real checkpoint, so the interface shows
recorded results rather than illustrative ones. The dissertation is generated by
script from the registered artifacts, so the document and the evidence cannot drift
apart.

What is still not built is the live inference service. There is no API, no
server-side model loading, and no deployment. The prototype demonstrates the
interface and the recorded results.

\section{Dataset and experimental setup}

\begin{table}[H]
\centering
\caption{PlantDoc label mapping}
\label{tab:mapping-phase3}
\tabfont
\begin{tabularx}{\textwidth}{@{}lY@{}}
\toprule
\textbf{Decision} & \textbf{Count} \\
\midrule
PlantDoc classes & 28 \\
Mapped to a PlantVillage target & 28 (25 confirmed, 3 confirmed with caveat) \\
Unmapped and excluded & 0 \\
Relationship recorded as same disease & 12 \\
Relationship recorded as healthy class & 10 \\
Relationship recorded as synonym or likely synonym & 5 \\
Relationship recorded as assumed synonym & 1 \\
Test images scored & 236 across 27 classes \\
\bottomrule
\end{tabularx}
\end{table}

The mapping was completed at mid-term and is unchanged: every PlantDoc class has a
target and a recorded relationship, and the three caveated mappings stay in the scored
set rather than being dropped, because dropping them would bias the field estimate
towards classes that happen to be easy to map.

The background-replacement augmentation required a corpus of field images that does
not overlap either evaluation set. A hash-disjointness gate runs before training and
refuses to start on any byte-identical overlap with PlantDoc or PlantVillage, so a
contaminated corpus cannot produce a result.

\section{Testing and validation}

The suite grew from 17 to 39 tests. Additions cover the calibration and reliability
path, data loading and split resolution, Grad-CAM output, and the checkpoint loader.

Three defects were found and fixed in this phase, each caught by checking output
rather than by checking exit status.

\begin{enumerate}[leftmargin=1.4em,itemsep=3pt]
\item A generated notebook ran its disjointness cell before the cell that writes the
      scripts, so the check could never have run. The builder now refuses to emit a
      notebook where a cell executes a script written by a later cell.
\item A field evaluation failed on the remote host with
      \texttt{ModuleNotFoundError} because the script resolved its project root one
      level above the working directory. Every subprocess now carries an explicit
      \texttt{PYTHONPATH}.
\item The first reliability figure connected bins with a line, which gave a bin
      holding one image the same visual weight as the bin holding 10{,}590 of
      10{,}709. It read as severe miscalibration that does not exist. The figure now
      encodes bin count as marker area and draws no connecting line.
\end{enumerate}

\section{Quantitative results and performance analysis}

Two model cohorts appear below and their field numbers differ. They are not
interchangeable: cohort A is the locked-test reference model used for the domain gap,
and cohort B is the paired augmentation experiment.

\begin{table}[H]
\centering
\caption{Domain gap, cohort A}
\label{tab:gap-phase3}
\tabfont
\begin{tabularx}{\textwidth}{@{}lrrrY@{}}
\toprule
\textbf{Measure} & \textbf{PlantVillage} & \textbf{PlantDoc} & \textbf{Change} & \\
\midrule
Accuracy, published checkpoint & 0.9928 & 0.1737 & $-81.9$ pp & \\
Accuracy, five-seed mean & 0.9928 & 0.2085 & $-78.4$ pp & \\
Macro-F1, published checkpoint & 0.9895 & 0.1798 & $-81.0$ pp & \\
\bottomrule
\end{tabularx}
\end{table}

PlantVillage figures are validation-set values for the five-seed reference model. The
published checkpoint and the five-seed mean are different quantities, so both are
shown.

On the locked test the reference model scores macro-F1 0.9931
(95\,\% interval $0.9910$ to $0.9951$), with expected calibration error 0.0022.

\begin{table}[H]
\centering
\caption{Intervention experiment, 298 validation images}
\label{tab:swap-phase3}
\tabfont
\begin{tabularx}{\textwidth}{@{}lrrY@{}}
\toprule
\textbf{Measurement} & \textbf{Value} & \textbf{95\,\% interval} & \\
\midrule
Chance adoption (1 in 38) & 0.0263 & --- & \\
Same-class control, flip rate & 0.0436 & 0.023--0.067 & the edit alone \\
Leaf swap adoption & \textbf{0.5034} & 0.446--0.557 & \\
Background swap adoption & \textbf{0.2181} & 0.171--0.268 & \\
Background minus leaf & $-0.2852$ & $-0.362$ to $-0.205$ & \\
\bottomrule
\end{tabularx}
\end{table}

The background carries real signal, at eight times chance, but the leaf carries more
than twice as much. This supersedes the mid-term statement that lab accuracy was
substantially background dependent, which was based on neutralisation alone and could
not distinguish dependence from disruption.

\begin{figure}[H]
\centering
\includegraphics[width=0.78\textwidth]{fig_swap.png}
\caption{How often the prediction follows the swapped-in part. Swapping the leaf
(dark blue) changes the prediction half the time; swapping the background (orange)
changes it about one time in five. Bars show 95\,\% intervals.}
\label{fig:swap}
\end{figure}

\begin{table}[H]
\centering
\caption{Paired augmentation experiment, cohort B}
\label{tab:h3-phase3}
\tabfont
\begin{tabularx}{\textwidth}{@{}lYYY@{}}
\toprule
\textbf{Condition} & \textbf{In-domain macro-F1} & \textbf{Field macro-F1} & \textbf{Per-seed field} \\
\midrule
Baseline & 0.9883 & \textbf{0.2232} & 0.2331, 0.2130, 0.2235 \\
Composite & 0.9796 & \textbf{0.2356} & 0.2287, 0.2451, 0.2330 \\
\midrule
Difference & $-0.0087$ & $+0.0124$ & $-0.0044$, $+0.0320$, $+0.0096$ \\
\bottomrule
\end{tabularx}
\end{table}

The paired difference is $+1.24$ points with a 95\,\% interval of $[-3.32, +5.80]$,
two of three seeds improved, and 1.6\,\% of the gap is recovered. There is no
detectable effect.

Calibration on field data rises from 0.0026 to 0.5808. At 90\,\% coverage the
reference model gains 0.08 points, so abstention does not produce a usable operating
point.

\begin{figure}[H]
\centering
\includegraphics[width=0.82\textwidth]{fig_coverage.png}
\caption{Accuracy on the field images the model is most sure about. If its confidence
could be trusted, each line would rise steeply as less-sure images are dropped
(moving left). They barely move; the SmallCNN line falls.}
\label{fig:coverage}
\end{figure}

\section{Integration and verification}

The registry now backs 23 claims, each tied to a run and an artifact directory.
\path{build_tables.py} with \texttt{-{}-verify} compares four generated tables against the registry
and fails on any mismatch, including a claim citing an unknown run. This is exercised
on every build.

Notebook generation verifies eight embedded files byte for byte against the repository
before producing the notebook, and refuses on any difference.

Cross-run consistency is checked by dataset fingerprint: all sweep runs report the
same split hash, and each metrics file records the learning rate actually applied per
epoch, so a schedule defect is visible in the artifact rather than inferred.

The final report is generated from the same artifacts, and the figures are generated
by script with no hand-entered values.

\section{Technical challenges and limitations}

The expensive failure of this phase was not computational. A notebook passed its
compile check and failed on first real execution because of cell ordering, and
because the error handler reported the failure as dataset contamination. That message
would have sent the work looking in the wrong direction. The lesson is recorded in the
build guard: a check that cannot distinguish a missing input from a failed test will
eventually report the wrong one.

A second class of error appeared repeatedly and is worth stating because it affects
how the results should be read. Several builds reported success while the artifact was
wrong, because a non-zero exit is not the only way to fail. Each was caught by opening
the output.

Open limitations. The field evaluation rests on 236 images from a single source with
four to twelve images per class, so the interval on a single accuracy estimate is
roughly ten points wide. Part of the gap is likely annotation noise, which this work
does not estimate. Three mappings are caveated rather than confirmed. The augmentation
result has three seeds, so the paired interval has two degrees of freedom and is wide.

\section{Milestone achievement and future work}

\begin{table}[H]
\centering
\caption{Milestone status at Phase III}
\label{tab:milestones}
\tabfont
\begin{tabularx}{\textwidth}{@{}lYl@{}}
\toprule
\textbf{Item} & \textbf{Deliverable} & \textbf{Status} \\
\midrule
H1 & Domain gap measured with intervals, five seeds & Complete \\
H2 & Mechanism measured by intervention & Complete \\
H3 & Background augmentation tested, paired, three seeds & Complete, null result \\
H4 & Calibration and abstention bounds & Complete \\
Splits & Frozen split with disjointness gate & Complete \\
Engine & Four architectures, one protocol, 14 runs, zero failures & Complete \\
Mapping & PlantDoc labels, 28 of 28 with caveats recorded & Complete \\
Tests & 39 tests plus six reproducibility gates & Complete \\
Registry & 23 claims gated against artifacts & Complete \\
Dissertation & Generated by script from registered artifacts; writing in progress & In progress \\
Explainability & Grad-CAM overlays for 18 field cases done; 50 hand-labelled maps pending & Partial \\
Front end & Static interface, recorded output, no live service & Partial by design \\
Field set & Second independent dataset with expert labels & Not started \\
OOD detection & Replacement for the confidence threshold & Not started \\
\bottomrule
\end{tabularx}
\end{table}

One point of consistency with the approved abstract belongs here rather than in a
reader's discovery. The abstract commits to a confidence threshold that flags
uncertain cases. The mechanism was implemented and measured. On field data it gained
0.08 points at 90\,\% coverage, so the expectation was not confirmed, and it is
reported as a negative finding rather than as a delivered feature.

The abstract also commits to targeted augmentation for realistic variation. That was
implemented and tested at three seeds per condition, and it did not move the field
gap. Both are reported as measured outcomes.

Remaining work is decision work and writing: obtain a second field dataset, decide
whether to pursue the Out-of-distribution direction, and agree the abstract wording
with the guide where the evidence no longer supports the original expectation.
